% !TEX root = ProjectDescription.tex
\subsection{Thrust 3: \thrustthree}
\label{sec:thrust3}

\noindent\textbf{Problem Statement.}
% TODO(glitching): 2-4 sentences.
%   - Existing fault-injection countermeasures are applied blanket (2-3x overhead) or by hand.
%   - Since the compiler creates the surface (Thrust 1) and we can locate the exploitable subset
%     (Thrust 2), the compiler can also shrink it -- selectively, at minimal cost, with guarantees.
%   - We reduce, not eliminate: full elimination is impossible under physical fault injection.
%     The defensible claim is provable resistance to a BOUNDED number of faults (fault model).
%   - Challenge: cost-optimal placement + a defensible claim of residual security.

\smallskip
\noindent\textbf{Research Questions.}
\begin{itemize}[leftmargin=*,nosep]
  \item \textbf{RQ3.1:} How can countermeasures be placed \emph{selectively} to reduce the exploitable surface (from Thrust~2) at minimal performance, size, and energy cost?
  \item \textbf{RQ3.2:} Can the compiler \emph{avoid introducing} the surface during code generation, rather than patching it afterward?
  \item \textbf{RQ3.3:} How can residual resistance be \emph{formally verified} under a bounded fault model, turning hardening into a guarantee rather than a heuristic?
\end{itemize}

\smallskip
\noindent\textbf{Approach Overview.}
% TODO(glitching): one paragraph tying the tasks below to the RQs above.

\subsubsection{\researchthreeone}
Uses the exploitable-site map from Thrust~2 to place fault-injection countermeasures selectively, minimizing performance, code-size, and energy overhead while maximizing surface reduction.
\noindent\textbf{Goal.} % TODO(glitching)
\noindent\textbf{Approach.} % TODO(glitching): vulnerability-guided, cost-aware countermeasure selection/placement.
\noindent\textbf{Expected Outcome.} % TODO(glitching)

\subsubsection{\researchthreetwo}
Modifies compiler instruction selection and register allocation to avoid generating code patterns that create single points of failure, preventing the compiler from introducing surface in the first place.
\noindent\textbf{Goal.} % TODO(glitching)
\noindent\textbf{Approach.} % TODO(glitching): fault-aware instruction selection / register allocation that does not create single points of failure.
\noindent\textbf{Expected Outcome.} % TODO(glitching)

\subsubsection{\researchthreethree}
Formally verifies that the hardened binary provably resists up to $N$ faults under the Thrust~2 fault model, turning hardening into a machine-checked guarantee rather than a heuristic.
\noindent\textbf{Goal.} % TODO(glitching)
\noindent\textbf{Approach.} % TODO(glitching): verify the hardened binary resists N faults under the Thrust 2 fault model.
\noindent\textbf{Expected Outcome.} % TODO(glitching)
